\documentclass[10pt]{article}
\usepackage{template}
\usepackage{amsmath,amssymb,microtype,tikz,needspace,xurl}
\usetikzlibrary{arrows.meta,positioning}
\setlength{\emergencystretch}{2.5em}
\SetKeys{date=2026-11-01,vol={2026-W44},title={每周重大新闻综述｜第44周},author={ChatGPT 编辑部},note={\small 观察期：2026年10月25—31日。截稿：11月1日上午（北京时间）。消息、政府声明、预测和预印本作者主张均按证据等级区分。}}
\newcommand{\weeklylead}[2]{\par\needspace{7\baselineskip}\bigskip{\Large\bfseries\color{newsblue}#1\par}{\footnotesize\color{newsgray}#2\par}\medskip\color{black}}
\newcommand{\weeklyfocus}[2]{\par\needspace{6\baselineskip}\bigskip{\large\bfseries\color{newsblue}#1\par}{\footnotesize\color{newsgray}#2\par}\smallskip\color{black}}
\newcommand{\weeklybrief}[2]{\par\needspace{5\baselineskip}\medskip{\normalsize\bfseries\color{newsblue}#1\par}{\footnotesize\color{newsgray}#2\par}\smallskip\color{black}}
\renewcommand{\refname}{参考资料}
\begin{document}
\begin{documents}
\raggedcolumns
\setlength{\parindent}{1.8em}
\setlength{\parskip}{2pt}
\weeklylead{霍尔木兹：一条航道，两套安全安排，和尚未落地的和平}{深度综述｜10月28—31日｜波斯湾外交、海上安全与能源}
10月的最后一周，波斯湾同时出现了两种看似相反的消息：商船通行有所恢复，地区各国却在为新一轮军事打击做准备。10月28日，《卫报》报道海湾国家正在推动分阶段协议，将制裁解除、霍尔木兹海峡开放、也门局势与地区安全安排放进同一个谈判框架。10月31日，路透社报道伊朗谈判人员与卡塔尔斡旋者在多哈继续会谈，前一天双方已在德黑兰进行了超过十二小时的接触。这里能确认的是外交活动确实在进行；不能确认的是美国与伊朗已经同意同一份文本，更不能将中间人的建议写成停战协议。\cite{iranReuters,iranGuardian}

谈判的障碍并不只是双方对某一项条件的价格判断不同，而是对顺序本身存在争议。伊朗希望解除港口封锁、释放被冻结的资产，并要求能够保障其航运利益的安排；美国方面强调核问题、核查与地区安全，认为不能在关键约束未得到保证时先解除压力。卡塔尔提出的短期建立信任方案据路透社报道未获伊朗接受。谈判若采取“先让航道开放，再讨论制裁”的路径，德黑兰担心让出筹码后难以取得承诺；若采取“先解除制裁，再开放航道”的路径，华盛顿及其伙伴又担心航运安全缺少可执行保障。和平方案因此需要写出同步履行、核验和违约补救，而不是只在标题上写下“交换条件”。

海峡的物理环境把外交争议放大了。它不是一条可以随意改道的陆上公路：油轮、液化天然气船、拖轮、港口调度、航行保险和船员安全共同决定实际运输能力。即便军事护航提高了通过数量，商业船东仍要评估船体损失、船员伤亡、货物交付和保单条款。10月30日《华盛顿邮报》报道，10月中旬原油出口量恢复到战前约八成，每日至少三十艘船舶通过，但仍有相当多的私人运营商回避该航道。这是报道所引用的估计，并非独立完整的海关或船舶逐艘清单。恢复到八成也不等于风险已经消失。\cite{hormuzWP}

价格传导有两个层次。第一层是实际供给：若船舶不能按时抵港，炼厂可能需要动用库存或改变采购。第二层是风险预期：即使尚未出现同等规模的现货短缺，市场也可能因为未来中断概率提高而抬高远期报价和保险成本。能源价格进一步影响化肥、货运、工业加热和居民出行，但不同国家受影响程度差异很大。进口国可能面对更高贸易账单，出口国也可能因为港口、保险与运输限制而无法完整获得价格上涨的收益。因此不能把油价上涨简单解释为“某一方获利”。

10月31日的局势进一步说明政治语言与实际行动必须分开。路透社报道，伊朗最高领袖要求军队为可能的袭击准备强力回应，美国国务卿鲁比奥则表示总统尚未决定是否采取新的军事行动，同时保留军事选项。伊朗官员关于核武政策可能重新讨论的说法，属于政治警告，不是核政策已经改变的正式文件；美国方面的威慑表态，也不能等同于新的打击命令已经下达。对于涉及核材料和核查机制的主张，还应等待国际原子能机构或其他可独立检验的材料。\cite{iranReuters,iranLeMonde}

更难的部分在于区域安全安排的参与者并不止美伊。沙特阿拉伯、卡塔尔、阿曼、土耳其和巴基斯坦各有不同的安全与经济利益，也门胡塞武装与红海航运的关系又使协议难以只局限于霍尔木兹。即便双方同意短期通航，未来谁监督航线、如何处置新发生的袭击、海上保险能否恢复正常、核查人员何时返回，仍然决定协议能否长期维持。海湾国家推动包含伊朗的安全机制，与美国希望维持的地区军事存在之间，也可能存在制度性张力。

另一个常被忽略的风险是合同的时间尺度。国际石油贸易的交货日期、租船合同、长期天然气供应协议和保险续保周期并不一致。某个星期的护航成果，可能无法立刻改变未来数月的保费与信用条件；反过来，即使本周没有新的重大损毁，企业也可能因为担心下次航运中断而提前增加库存。评价外交谈判的经济成果，需要观察价格、运量和合同条款是否在较长时间内共同改善，而不是只截取一个交易日的油价变化。

新闻报道应避免把任何一个政府的立场当作中立事实。伊朗方面有动机强调遭受封锁的经济损失，美国政府有动机突出其护航与威慑成果，海湾国家希望避免本国基础设施再次成为战场，航运企业则关注可投保性和可预测性。《卫报》侧重外交方案，《华盛顿邮报》侧重军事护航的运行情况，路透社追踪谈判和公开表态；它们各自提供了不同类型的证据，而不是三份彼此完全独立的现场核验。

截至10月31日，本刊可以确认外交接触与军事准备并存，无法确认最终协议、海峡全面安全开放或战争即将重启。接下来最有价值的观察对象不是强硬措辞本身，而是协议草案、航行通告、核查安排、船舶实际通过量和保险市场是否持续恢复。和平不是一句宣布，而是一套在冲突双方都不信任对方时仍能执行的程序。

\begin{center}
\begin{tikzpicture}[>=Stealth,font=\scriptsize,node distance=2mm,every node/.style={align=center}]
\node[draw=newsblue,rounded corners,fill=blue!5,text width=2.35cm,minimum height=0.8cm] (a) {海上安全风险\\袭击、封锁与护航};
\node[draw=orange!70!black,rounded corners,fill=orange!6,text width=2.35cm,minimum height=0.8cm,right=3mm of a] (b) {保险与有效运力\\等待、航线、成本};
\node[draw=newsblue,rounded corners,fill=blue!5,text width=2.35cm,minimum height=0.8cm,below=5mm of a] (c) {供给和交付预期\\库存与风险溢价};
\node[draw=orange!70!black,rounded corners,fill=orange!6,text width=2.35cm,minimum height=0.8cm,right=3mm of c] (d) {终端能源成本\\运输、工业、居民};
\draw[->,thick] (a)--(b);\draw[->,thick] (a)--(c);\draw[->,thick] (b)--(d);\draw[->,thick] (c)--(d);
\end{tikzpicture}

{\footnotesize\color{newsgray}图1｜编辑部原创传导机制示意（2026-11-01）；非实测数据，不表示某次袭击已经独立核实。}
\end{center}

\weeklylead{乌克兰局部停火：救援成功，为什么还不是全面和平？}{深度综述｜10月27—31日｜奥列什基撤离、俄乌会谈与防空援助}
乌克兰南部被俄罗斯占领的奥列什基，在本周提供了战争新闻中罕见的一条具体进展。10月27日起实施的局部停火为清理道路、运送救援物资和撤离平民创造了条件。10月31日，路透社报道部分居民经克里米亚、俄罗斯和白俄罗斯绕行，最终抵达乌克兰控制地区；他们描述了长期缺少食物、水、电力和医疗保障的生活。报道没有给出经过完整核验的最终撤离总人数，因此本刊不把尚在变化的数字写成最终成果。\cite{oleshkyReuters,oleshkyGuardian}

这项行动值得从人道法而非宣传胜负的角度理解。战争中的平民有权获得保护，救援需要交战方在时间、通道、安全检查、地雷处置和目的地选择上达成操作性安排。对年长者、病人和行动不便者而言，撤离并不是登上一辆车那么简单：长途转运、跨境文件、医疗转诊、家属联系和宠物安置都可能决定一个人是否愿意离开。局部停火的实际价值，首先体现在这些具体生活条件的改善。

局部停火却不能直接推广为全国停火。一个地区短暂停止军事行动，可以依赖有限范围内的协调；全国范围的停止射击则涉及漫长接触线、远程导弹与无人机、边境空域、海上行动和军事补给。监督机构需要知道什么构成违约、由谁取证、如何回应意外冲突，以及交战方是否允许第三方进入。若这些问题悬而未决，双方都可能担心对方利用停火重新部署。人道通道的成功证明有限合作有可能，并不自动解决领土和安全保证争端。

本周俄罗斯与乌克兰仍在讨论美国斡旋的下一轮三方会谈。乌克兰方面希望11月6日至10日前后在阿布扎比举行会议，并希望欧洲参与；俄罗斯外长拉夫罗夫10月30日在G20外长会议期间表示俄方立场未变，反对乌克兰成为俄罗斯所称的安全威胁。双方对领土、军力和长期安全保障的要求仍然相距甚远。会谈时间的公开建议不等于日程已经由三方共同确认，更不等于已经存在可执行的和平文本。\cite{lavrovReuters,ukraineTalks}

另一个容易引起误读的消息涉及乌克兰对俄罗斯炼油设施的远程打击。10月30日，路透社转述《华尔街日报》引述匿名消息人士的报道，称乌克兰在美国压力下暂时暂停部分行动；路透社明确表示无法独立核实。匿名报道可以作为需要追踪的线索，却不能与乌克兰政府发布的正式命令或实际军事行动记录等同。相关说法涉及美国国内能源价格、俄方石油收入以及乌克兰如何回应本国能源设施遭袭，任何单一解释都过于简单。\cite{refineryReuters}

美国同时批准一项拟议中的爱国者防空拦截弹销售，媒体报道涉及约五百枚拦截弹和二十一亿美元规模。这是对外军售的审批进展，不是五百枚导弹已经运抵乌克兰。实际效果还取决于资金来源、合同签署、生产能力、交付时序、系统适配和训练。防空导弹的战略意义，也不能只以单枚价格衡量：它涉及保护人口密集区、能源基础设施和军事设施之间的艰难分配。\cite{g20AP}

从制度设计上看，人道行动还需要回答撤离后的问题。离开占领区的居民是否有稳定住处、医疗保险、身份证明和家庭团聚渠道？如果房屋损毁或财产记录丢失，未来重返故乡时如何主张权利？临时停火可以打开通道，却不能自动建立长期安置体系。国际人道组织和地方政府在救援结束后的持续协调，可能比一次媒体关注的跨境旅程更耗时。对选择留下的居民，也需要有可持续的补给与医疗方案；不应把“没有撤离”简单解释为不存在危险。

局部人道行动的功劳归属也应谨慎。白宫强调美国第一夫人及其他斡旋方的作用，乌方则指出相关协调此前已经进行。不同参与者可能都在某个环节发挥作用，但本刊没有完整的谈判记录来精确分配个人贡献。相比谁在新闻稿中得到最多赞誉，更重要的是居民是否安全抵达、是否能够自由选择目的地，以及剩余居民能否持续获得援助。

与之相关的另一层难题是信息的可验证性。交战双方都掌握部分现场资料，平民证言能够揭示生活处境，却未必能够核实全部军事行动；卫星影像可以显示建筑与道路变化，却不能单独说明谁在何时作出了特定命令。新闻机构应将现场采访、公开影像、行政记录和国际组织的后续报告相互对照，既避免因证据不足而抹去真实苦难，也避免把未经确认的战场数字当成既定历史。

这周的现实是两条线并行：一个地方的居民得到了离开危险环境的机会，其他地区的战争并未结束。把局部停火称为没有意义，会低估它对生命的价值；把它写成全面和平的前奏，则会高估现有证据。下一周应关注撤离总人数的独立核实、持续补给、停火是否延长、三方会议是否举行，以及停火监督与战后安全安排是否出现可供检验的书面文本。

\weeklylead{欧洲央行按兵不动：能源通胀、债务压力与行长提前离任}{深度综述｜10月29日｜货币政策、财政与制度独立性}
10月29日，欧洲中央银行管理委员会决定维持三项关键利率不变。其官方声明把困难说得相当清楚：能源价格持续上涨，增加通胀上行风险，同时压低经济增长前景；融资条件自上次会议以来已经收紧，又会对需求和价格形成抑制。与此同时，行长克里斯蒂娜·拉加德宣布将在2027年2月7日提前离任，较原定任期结束时间早约九个月。这两件事在同一天发生，使市场必须同时判断政策路径与机构领导层的连续性。\cite{ecbOfficial,lagardeReuters,lagardeAP}

理解这次决定，首先要区分需求冲击与供给冲击。若通胀主要来自居民和企业支出过旺，提高利率可以通过借贷成本抑制需求；但若油气运输受阻导致能源成本上涨，提高利率并不能直接制造更多天然气或让受损航道恢复安全。央行真正能够影响的是二次传导：企业是否持续涨价，工人是否要求补偿性加薪，公众是否相信未来通胀仍会回到目标。货币政策因此要在控制预期和避免过度压制增长之间权衡。

欧洲央行声明指出，基础通胀目前总体稳定，工资尚未出现显著反应；但能源冲击的幅度和持续时间可能使后续风险上升。这解释了为何它既没有立即加息，也没有承诺未来不再加息。政策声明强调逐次会议依据数据决定，而非预先绑定某条利率路径。市场将这种表述解读为偏向谨慎，不能把交易日的汇率变化当作央行已秘密决定12月政策的证据。\cite{ecbOfficial,ecbReuters}

利率传导并不是一条机械直线。政策利率影响银行资金成本和市场预期，随后通过抵押贷款、企业融资、政府债券收益率和汇率改变经济行为。企业面对高能源成本和高融资成本时，可能推迟投资；家庭则可能在房贷与取暖费用之间被迫重新分配预算。不同欧元区成员国的财政状况和银行结构不同，政策收紧产生的影响也不均匀。

法国等高负债国家的财政处境尤其值得关注。政府需要为旧债再融资，市场收益率上升可能增加未来利息支出；如果同时削减公共服务或增加税负，又可能引发社会反弹。央行必须维持货币稳定，却不能替代议会决定税收和支出。欧元区的制度设计要求把共同货币政策与各国财政政策区分开，这在能源冲击下变得更加棘手。

拉加德提前离任的解释涉及制度边界。她表示希望能够更自由地参与货币政策以外的公共讨论，认为现任央行官员受中立性要求约束。外界也注意到她离任时间早于法国总统选举，可能影响继任者任命的政治节奏；但“时间上有利于某种任命安排”与“提前离任就是为了影响选举”不是同一个可以直接证明的命题。对动机的分析应标明推测性质。\cite{lagardeReuters,lagardeAP}

还有一个经济学上的时间差问题：居民今天看到的是加油站和超市的价格，央行能够改变的却主要是未来数季的融资与需求。利率调整需要通过贷款合同重定价、企业投资和工资谈判才能逐渐传导。若政府在这段时间实施普遍能源补贴，短期生活成本可能缓和，但财政赤字与能源消费激励也会改变；若采取面向低收入家庭的定向支持，则需要行政能力和公众信任。央行和政府可以有不同的职责，但不能假装彼此政策互不影响。

央行独立性不是要求其成员没有个人意见，而是要求利率决策和资产购买等政策不能轻易受短期选举压力支配。行长的公开讲话本身又会影响市场，因此在任时的表达受到额外约束。继任者无论来自哪个国家，都需要面对相同的制度难题：如何在能源冲击、公共债务和政治压力之间维护通胀目标，同时向公众解释政策为何可能增加短期痛苦。

制度上还存在欧元区内部差异。某些国家家庭持有浮动利率贷款的比例较高，政策利率变化较快进入家庭预算；另一些国家长期固定利率贷款占比更高，冲击传导较慢。能源进口依赖、工业结构和财政空间也不同。统一利率只能在全区层面作出权衡，不能保证每个成员国都恰好得到最适合自身的政策。这也是为何债券市场分化与成员国财政政策始终是欧洲央行必须观察的背景变量。

这周最值得关注的不是欧洲央行是否给市场一个令人满意的方向。而是它拒绝在证据尚不充分时作出确定承诺。下一步应观察油气价格、工资谈判、通胀预期、主权债利差和12月会议的新数据；继任人选则应由正式任命程序核实。能源冲击可以在几周内改变价格，却可能用更长时间改变公众对货币和政府的信任。

\weeklyfocus{G20 外长会：没有联合公报，是否就没有外交成果？}{国际外交｜10月30—31日｜路透社、美联社}
10月30日至31日，二十国集团外长在美国亚特兰大举行会议。美联社报道，会议结束时未能就关键安全议题形成联合公报，而是以主席声明总结讨论。俄罗斯外长拉夫罗夫受邀出席，引起部分欧洲国家反对，传统集体合影因此取消。会议同时涉及乌克兰战争、中东局势、能源与粮食安全、海上通行及关键矿产供应链。\cite{g20AP,g20Reuters}

联合公报需要较广泛共识，主席声明则主要说明主席方如何概括会议。两者的政治分量和责任归属不同；未能达成联合公报确实反映成员分歧，却不意味着会场内所有双边接触、议程协调和信息交换都没有价值。反过来，会议顺利召开也不能证明参与者已经就战争或贸易争议形成解决方案。

美国国务卿鲁比奥认为和平需要与冲突各方对话，欧洲外交官则质疑在俄军持续攻击乌克兰时给予俄方高级代表正常外交待遇的象征意义。俄方强调其所称的安全利益，乌克兰并未参加这次G20外长会议。这些立场之间存在实质冲突，不能用“双方各执一词”就抹平国际法、战争责任与平民保护问题。

本周还举行了美、日、印、澳四方安全对话外长会议。四国声明强调供应链韧性、关键矿产和海上通行；其对所谓经济胁迫的担忧虽未直接点名中国，却显然处于中美技术与贸易竞争背景中。判断这些倡议是否真正产生效果，应看项目资金、执行机构和后续协议，而不只是共同声明中的形容词。\cite{quadReuters}

本刊认为，G20缺少联合公报不是制度死亡证明，而是多边治理在地缘对立下的可见成本。若主要国家连对战争性质和航运安全的最低表述都难以达成一致，未来协调能源危机、债务风险与供应链冲击也会更加困难。

\weeklyfocus{中国十月制造业再度收缩：PMI 的 49.3 究竟说明什么？}{经济观察｜10月31日｜官方调查与独立报道}
10月31日公布的中国官方制造业采购经理指数（PMI）为49.3，低于9月的50.1。美联社报道，新出口订单指数为48.7，生产指数为49.8；《华尔街日报》还指出非制造业商务活动指数降至49.7。采购经理指数以50为扩张和收缩的分界，49.3意味着受调查企业的经营状况总体呈收缩信号，而不是制造业总产值当月下降了0.7%。\cite{chinaPMIAP,chinaPMIWSJ}

PMI是一种扩散指数。调查人员统计企业报告某项活动改善、持平或恶化的比例，再计算指标；它主要反映变化方向与广度，而非产值或利润的精确变化幅度。某些行业可能继续增长，另一些行业则明显走弱。把49.3解读为“所有工厂都在衰退”同样不准确。

这次数据值得结合内需与外需分别观察。新出口订单下降说明外部订单环境承压，但不能仅凭这一项就确定是哪个国家的关税、汇率或需求变化造成的；国内房地产调整、居民消费谨慎、企业投资预期和库存周期也会影响制造业。官方指数的样本结构、节假日效应和季节调整都需要留意。

中国出口竞争力与国内消费疲弱可以同时存在。若企业在国内需求不足时增加出口，短期可能维持生产，却会在海外市场引起更强的贸易摩擦。美联社报道指出欧洲等贸易伙伴对贸易不平衡表示担忧。这些争议需要分别核对贸易额、产品类别和双边统计口径，不能把一个月的PMI当作国际贸易政策的完整因果解释。

接下来应与工业增加值、出口实绩、价格指数、企业利润和居民收入交叉观察。PMI为经济放缓提供了及时但不完整的信号；真正的政策问题仍是如何改善内需、恢复企业投资信心，以及避免过度依赖某一出口行业。

\weeklyfocus{NASA 的 378 天模拟火星任务：封闭舱里研究的不是火星风景}{航天与生命科学｜10月27—31日｜NASA CHAPEA 第二次任务}
美国国家航空航天局在10月27日公告，第二次“乘员健康与表现探索模拟”（CHAPEA）任务的四名志愿者计划于10月31日结束378天隔离生活。四人自2025年10月19日起居住在休斯敦约1700平方英尺的3D打印模拟栖息地中，执行作物种植、设备维护、模拟舱外活动与资源管理等任务。NASA的预告可确认任务设计与预定结束日期；本刊没有取得足以独立核实全部后续医学数据的材料。\cite{nasaCHAPEA}

模拟火星生活的核心难题并非制造一张像火星的背景图，而是在人与地面控制存在通信延迟、补给受限、活动空间有限的情况下保持长期工作能力。飞船任务还必须面对设备故障、睡眠与昼夜节律、团队冲突、食品结构以及紧急医疗决策。地球上的隔离舱不能模拟火星重力或宇宙辐射，但能对组织流程和行为因素进行受控研究。

研究者需要比较的是长期行为变化，而不是仅凭离舱时参与者的精神状态下结论。心理压力、认知表现、睡眠、团队协作、营养和生理指标可能随任务阶段变化，且样本只有四人，个体差异非常大。若某项表现改善，还需要判断是训练效果、任务适应还是测量条件变化。

CHAPEA也提醒人们，深空任务的系统工程并不只属于推进器与轨道力学。食物、医疗、通信带宽和乘员时间都是有限资源；提高一个环节的冗余，往往会占用另一个环节的质量、体积或能源预算。NASA表示数据将帮助未来月球和深空任务规划，但数据采集结束并不等于所有结论已经发表。

\weeklyfocus{数学拾趣｜从格点间望出去：六维与七维的视线遮挡}{数论／组合／几何｜10月26日 arXiv:2610.37711｜Vineet Kahlon}
想象平面上每个半整数格点都放着一个小圆盘，观察者站在原点，朝第一象限任意方向发射一条射线。圆盘半径至少多大，才能保证没有一条射线从所有圆盘之间的缝隙穿过？把圆盘换成高维欧氏空间中的球，就得到视线遮挡问题的一种形式。问题看起来像几何覆盖，却很快牵涉整数方向向量、格的对偶与丢番图逼近。\cite{mathView}

10月26日，Vineet Kahlon在arXiv提交预印本，声称精确计算了第六维和第七维的对应极值：第六维为107/312，第七维为73/180。这里的数值对应论文定义的平方距离极值量，而不是未经开方就能直接作为球半径的数值。作者还给出取得极值的整数方向向量，并称可通过坐标置换描述全部极值情形。

为何高维问题能够变成有限搜索？作者摘要称，证明利用有理子空间上的精确下界、匹配的纤维上界以及Banaszczyk转移定理，并沿着对偶格的Korkine--Zolotarev旗标逐步约束候选。搜索过程使用精确有理数运算，作者称每个分支均已解决，构成可核查的计算证书。这与单纯在大量随机方向上测试并猜测最优值不同。

读者应注意两个证据层次。首先，本文确实存在，作者、日期和摘要可从arXiv核对；其次，极值定理及计算证书的正确性仍需逐项检查。预印本不等于同行评议完成，本刊没有复现搜索程序，也没有独立证明上述高维结论。这个问题的美感在于，一个“能否看穿圆球阵列”的直观图景，最终与非常精细的格几何发生联系。

\weeklyfocus{数学拾趣｜黎曼猜想可以被约数和不等式检测吗？}{解析数论／计算实验｜10月29日 arXiv:2610.39120}
对正整数n，记其所有正约数之和为sigma(n)。例如12的正约数为1、2、3、4、6、12，因此sigma(12)=28。Robin判别准则把著名的黎曼猜想与约数和增长速度联系起来：对于所有大于5040的整数，某个涉及sigma(n)、欧拉常数和log log n的不等式成立，当且仅当黎曼猜想成立。一个关于复变函数零点的命题，竟然能改写成纯整数不等式，这正是解析数论的魅力。\cite{mathCA2}

10月29日，Dmitry R. Leshchiner、Oleg R. Musin与James Thompson提交预印本，研究所谓“极大丰盛数”（colossally abundant numbers）及其更小的CA2子类。作者在摘要中指出，已有理论可以将相关判别测试缩减到特殊整数集合；他们进一步研究是否还能减少需要检查的候选，并对Ramanujan条件性界给出的衰减规律开展数值比较。

为什么选择特殊整数而不是逐个检查所有n？因为在一些优化问题中，最容易让不等式接近失效的对象具有高度特殊的质因数结构。研究这些极端候选，可以让无限整数集合上的问题具有更清晰的组织方式。但必须强调：计算到再大的有限范围，也不能直接证明一个要求对所有整数成立的命题。

论文讨论归一化差值、黎曼zeta函数零点引起的振荡表达式，以及截断零点求和造成的数值误差。作者报告计算范围内仍存在缓慢变化的部分，并提出关于收敛和高阶项的问题。这是对判别准则的数值和理论研究，不是宣布黎曼猜想已被证明。作为10月29日的arXiv预印本，它尚不能因出现“黎曼猜想”字样就获得已被同行核验的地位；本刊也未独立检查其计算代码。

\weeklybrief{四方安全对话讨论供应链与可信 AI}{国际科技政策｜10月31日}
美、日、印、澳四国外长在G20会议间隙重申供应链韧性与海上通行，讨论关键矿产、半导体供应链风险及区域可信人工智能项目。路透社报道，四方宣布涉及八千五百万美元以上的相关倡议；这属于承诺与计划，不能当作资金已经全部支出或项目已实现效果。\cite{quadReuters}

\weeklybrief{本周未刊登的数学与文化线索}{编辑说明｜10月25—31日}
本刊检索了数论、组合、几何及Lean/mathlib相关预印本，并选取两篇可从arXiv直接核实元数据的论文详述；没有为凑栏目而给未经检查的数学主张加上“重大突破”标签。文化与其他科学方向同样经过检索，但本期可核实材料与篇幅优先用于战争、人道救援、宏观政策和数学研究。未刊登不代表这些领域本周没有进展。

\par\medskip
\begingroup\footnotesize
\bibliographystyle{plain}
\bibliography{refs}
\endgroup
\end{documents}
\end{document}
